% Generated draft: requires mathematical and source review.
\documentclass[11pt]{article}
\usepackage[margin=25mm]{geometry}
\usepackage{fontspec}
\setmainfont{Latin Modern Roman}
\usepackage{amsmath,amssymb,mathtools,unicode-math}
\setmathfont{Latin Modern Math}
\usepackage{enumitem,needspace,xurl,hyperref}
\setlength{\parindent}{0pt}
\setlength{\parskip}{5pt}
\setlength{\emergencystretch}{3em}
\newcommand{\N}{\mathbb{N}}
\newcommand{\Z}{\mathbb{Z}}
\newcommand{\Q}{\mathbb{Q}}
\newcommand{\R}{\mathbb{R}}
\newcommand{\C}{\mathbb{C}}
\newcommand{\F}{\mathbb{F}}
\newcommand{\A}{\mathbb{A}}
\newcommand{\PP}{\mathbb{P}}
\newcommand{\E}{\mathbb{E}}
\newcommand{\eps}{\varepsilon}
\newcommand{\dd}{\,\mathrm{d}}
\begin{document}
\section*{UnsolvedMath ID: 10400183}
\textbf{Source problem code:} catalogue entry
\subsection{Definitions and mathematical statement}
The conjecture is that the integral diagram group \(\mathcal A(\varnothing;\mathbb Z)\) is torsion free. In the usual closed-Jacobi-diagram interpretation, start with the free abelian group generated by finite trivalent graphs with cyclic orderings at vertices and no external skeleton or univalent legs, then quotient by antisymmetry and IHX relations. Disconnected and empty diagrams, allowed graph degeneracies and any completion convention must match the original source.

The requested assertion is
\[
 \forall m\in\mathbb Z\setminus\{0\}\ \forall x\in\mathcal A(\varnothing;\mathbb Z):
 \quad mx=0\Longrightarrow x=0.
\]
Equivalently, the natural rationalization map
\[
 \mathcal A(\varnothing;\mathbb Z)\longrightarrow
 \mathcal A(\varnothing;\mathbb Z)\otimes_{\mathbb Z}\mathbb Q
\]
is injective. Passing to rational coefficients at the outset would erase exactly the torsion at issue and is not a replacement for the statement.

The excerpt does not reproduce the integral diagram presentation, loop/multiple-edge or saturation conventions, or grading/completion. These matter to the exact group. Its trailing text, ``10.2.2 Do finite type invariants distinguish homology 3-spheres?'', appears to be the following subsection heading rather than an additional assertion of Conjecture 10.8; no separate detection problem is silently appended to this ID.

The summary supplies only an open/triage description and no exact current theorem. The historical conjecture remains source information pending primary-literature review.
\subsection{Short English statement}
Does the integral closed-diagram group \(\mathcal A(\varnothing;\mathbb Z)\) have no torsion? The integral presentation must be fixed, and rationalizing first would lose the phenomenon being tested.
\subsection{Sources}
\begin{itemize}
\item Local registry: \url{data/unsolvedmath/problems.json}, permanent UnsolvedMath ID \(10400183\), source code catalogue entry; supplied statement, background, and research summary, updated 2026-08-06.
\item Ohtsuki, \emph{Problems on invariants of knots and 3-manifolds} (2002), Conjecture 10.8, PDF page 146: \url{https://msp.org/gtm/2002/04/gtm-2002-04-024s.pdf}. Supplied PDF-font-structure extraction accessed 2026-07-29; provenance distinguishes 231 numbered items.
\item No full preceding integral-diagram convention or separate resolution reference is supplied.
\end{itemize}
\end{document}
